% TOP_PROBLEM: {"id": 1347, "title": "Partition Principle Implies Axiom of Choice", "category_id": 10, "category": {"id": 10, "name": "set_theory", "display_name": "Set Theory", "description": "Foundations of mathematics, infinite sets, and cardinality.", "slug": "set-theory", "order_index": 10, "created_at": "2026-07-31T15:26:25.671Z"}, "status": "open"}
% FIRST_POSTED: 2026-10-01
% ATTEMPT_STATUS: unresolved
% Imported from: unsolvedmath_additional_500.tex; UnsolvedMath ID 1347
% !TeX program = lualatex
% Compile twice: lualatex 1347.tex
% Self-contained: no external bibliography, images, or \input files.
% Numeric section labels and visible numbers are original UnsolvedMath IDs.
\documentclass[11pt,a4paper]{article}
\usepackage[margin=24mm,headheight=14pt]{geometry}
\usepackage{fontspec}
\setmainfont{Latin Modern Roman}
\setsansfont{Latin Modern Sans}
\setmonofont{Latin Modern Mono}
\usepackage{amsmath,amssymb,mathtools}
\usepackage{unicode-math}
\setmathfont{Latin Modern Math}
\usepackage{microtype}
\usepackage{enumitem,xurl,needspace}
\usepackage[dvipsnames]{xcolor}
\usepackage{hyperref}
\hypersetup{unicode=true,colorlinks=true,linkcolor=MidnightBlue,urlcolor=MidnightBlue,
 pdftitle={UnsolvedMath Problem 1347},pdfauthor={Source-annotated compilation},bookmarksnumbered=true}
\usepackage{bookmark,tocloft,titlesec,fancyhdr}
\setlength{\cftsecnumwidth}{4.2em}
\cftsetpnumwidth{2.5em}
\cftsetrmarg{4em}
\titleformat{\section}{\Large\bfseries\raggedright}{\thesection}{0.7em}{}
\pagestyle{fancy}\fancyhf{}
\fancyhead[L]{\small\sffamily Additional UnsolvedMath statements}
\fancyhead[R]{\small Original catalogue IDs}
\fancyfoot[C]{\thepage}
\setlength{\parindent}{0pt}
\setlength{\parskip}{5pt plus 1pt}
\setlength{\emergencystretch}{3em}
\setcounter{tocdepth}{1}\setcounter{secnumdepth}{1}
\setlist[itemize]{leftmargin=3em,itemsep=4pt,topsep=4pt,parsep=0pt}
\allowdisplaybreaks
\newcommand{\N}{\mathbb{N}}
\newcommand{\Z}{\mathbb{Z}}
\newcommand{\Q}{\mathbb{Q}}
\newcommand{\R}{\mathbb{R}}
\newcommand{\C}{\mathbb{C}}
\newcommand{\F}{\mathbb{F}}
\newcommand{\A}{\mathbb{A}}
\newcommand{\PP}{\mathbb{P}}
\newcommand{\E}{\mathbb{E}}
\newcommand{\eps}{\varepsilon}
\newcommand{\dd}{\,\mathrm d}
\newcommand{\sref}[2]{\hyperlink{source:#1:#2}{[S#2]}}
\newif\ifshowcatalogue
\showcataloguetrue
\newcommand{\cataloguescope}[1]{\ifshowcatalogue
 \paragraph{Statement recorded in the supplied source.}
 #1\fi}
\begin{document}
% UnsolvedMath ID 1347; source catalogue code ST-001

\setcounter{section}{1346}

\section{Does the Partition Principle Imply Choice?}\label{1347}

{\small\sffamily UnsolvedMath ID 1347 · Set Theory}

\subsection{Definitions and mathematical statement}

\paragraph{Shared notation.}Unless a section says otherwise, $\N=\{1,2,\ldots\}$,
$\Z$ is the ring of integers, $\Q,\R,\C$ are the rational, real, and complex
fields, $[N]=\{1,\ldots,\lfloor N\rfloor\}$, and $\log$ is the natural
logarithm. For real functions of a parameter tending to infinity,
$f\ll g$ and $f=O(g)$ mean $|f|\leq Cg$ eventually for some fixed $C$;
$f\gg g$ reverses this comparison; $f=o(g)$ means $f/g\to0$;
$f\sim g$ means $f/g\to1$; and $f\asymp g$ means both order bounds.
A subscript on $O$ or $\ll$ specifies allowed constant dependence.
The expression $n^{o(1)}$ in an upper bound means growth smaller than
$n^\epsilon$ for each fixed $\epsilon>0$. Local definitions govern any
exception, finite-field convention, or use of the same symbol for another
object. An asymptotic claim is not automatically an effective algorithm.



Work in ZF set theory, without assuming the axiom of choice. The partition principle says that whenever there is a surjection $f:A\to B$, there is an injection $g:B\to A$. It does not require $f(g(b))=b$. The axiom of choice asserts a choice function for every family of nonempty sets. Ask whether ZF proves that the former principle implies the latter. \sref{1347}{1} \sref{1347}{2}

\paragraph{Statement recorded in the supplied source.}
Does the partition principle (PP) imply the axiom of choice (AC)? \sref{1347}{1}

\subsection{Short English statement}

Does the mere ability to inject the parts of every partition back into the original set already force the full axiom of choice?

\subsection{Research attempt: a choice-free consequence of PP and the failed section argument}
\textbf{Status and source correction.} All deductions below take place in ZF. We do not prove PP implies AC or construct a separating model. The inspected arXiv record for [S2] lists Adonai S. Sant'Anna, Otavio Bueno, Marcio P. P. de França and Renato Brodzinski. Its abstract claims a model through its Flow framework. We do not assume that claimed construction or certify it from the abstract; the arguments here are independent of it.

\subsubsection{1. The stronger statement really is equivalent to AC}
If every surjection $f:A\to B$ has a right inverse $s:B\to A$ with $f(s(b))=b$, apply this to the projection from the disjoint union $\bigsqcup_{i\in I}X_i$ of any family of nonempty sets onto $I$. The right inverse selects one member of each $X_i$, giving AC. Conversely AC chooses one member of each nonempty fiber of any surjection and gives a right inverse. A right inverse is injective, so AC implies PP.

PP supplies only an injection $j:B\to A$. The composite $f\circ j$ need not be onto, let alone the identity. Even in finite sets take $A=\{a_0,a_1,a_2\}$, $B=\{0,1\}$, let $f(a_0)=f(a_1)=0$, $f(a_2)=1$, and take $j(0)=a_0,j(1)=a_1$. Then $f\circ j$ is constant. This does not refute PP implies AC; it refutes the proposed one-step conversion of an arbitrary injection into a section by inverting its composite.

\subsubsection{2. PP rules out amorphous sets}
A set $A$ is amorphous when it is infinite but cannot be partitioned into two infinite subsets. We prove PP excludes such sets. Let $\operatorname{Fin}(A)$ be the set of finite subsets of an arbitrary infinite $A$. The function
\[\operatorname{Fin}(A)\longrightarrow\mathbb N_0,\qquad F\longmapsto|F|
\]
is surjective: ZF proves that an infinite set contains an $n$-element subset for each individual finite $n$. This statement does not simultaneously choose such subsets. PP now gives an injection $j:\mathbb N_0\to\operatorname{Fin}(A)$, so the finite sets $F_n=j(n)$ are all distinct.

Their union cannot be finite, since a finite set has only finitely many subsets. Recursively choose the least index $n_k$ for which $F_{n_k}$ contains a point outside the finite union of previously selected $F_{n_i}$, and put
\[B_k=F_{n_k}\setminus\bigcup_{i<k}F_{n_i}.\]
Such an index exists at every stage; the union of all $F_n$ is infinite while the previously selected union is finite. Each $B_k$ is nonempty and finite, and the blocks are pairwise disjoint. The recursion chooses least natural-number indices, not arbitrary elements of the blocks, so it uses no form of countable choice. Both $\bigcup_k B_{2k}$ and $\bigcup_k B_{2k+1}$ are infinite, because each contains infinitely many disjoint nonempty blocks. Put any remaining elements of $A$ into one side to get a partition of $A$ into two infinite subsets. Hence an amorphous set contradicts PP.

\subsubsection{3. A solvable case and the exact missing repair}
If the domain $A$ of a surjection is already well-ordered, its fibers have least elements, giving a right inverse in ZF without PP. For a general set no such fiberwise order is available. The block argument above also does not supply a point from each $B_k$: pairwise disjoint nonempty finite blocks need not come with chosen representatives in ZF. Replacing the least-index recursion by arbitrary element choices would hide exactly the kind of choice principle under investigation.

The adversarial audit therefore checks every selection in the proof: finite existence is used one size at a time, the single injection is supplied by PP, and all recursive choices are from well-ordered natural numbers. The result is the precise necessary consequence ``PP implies no amorphous sets.'' It is weaker than AC and provides no equivalence. A mechanism turning the unrestricted injection promised by PP into simultaneous fiber representatives, or a fully verified ZF model showing this impossible, remains missing.

\subsection{Sources}

{\small

\begin{itemize}

\item[\hypertarget{source:1347:1}{S1}] ULAM.AI, \emph{UnsolvedMath}, supplied \texttt{problems.json}, record 1347 (ST-001), statement and background. The embedded literature-review date is 2026-08-17; that is the archive’s review date, not a fresh certification. Dataset landing page: \url{https://huggingface.co/datasets/ulamai/UnsolvedMath}.

\item[\hypertarget{source:1347:2}{S2}] A. S. Sant'Anna, O. Bueno, M. P. P. de França and R. Brodzinski, Flow: the Axiom of Choice is independent from the Partition Principle, arXiv:2010.03664. \url{https://arxiv.org/abs/2010.03664}. Bibliographic information imported from the supplied record.

\item[\hypertarget{source:1347:3}{S3}] A Gentle Introduction to the Axiom of Choice, Mathematical Intelligencer (2026). \url{https://doi.org/10.1007/s00283-026-10531-4}. Bibliographic information imported from the supplied record.

\end{itemize}

}

\label{end:1347}
\end{document}
